\chapter{Conclusion}

\section{Overview}

This dissertation examined how prompt engineering influences correction behaviour in LLM-based grammatical error correction (GEC). The study was concerned not only with whether a prompt improves ERRANT accuracy, but also with whether it causes the model to change the learner's sentence more than necessary.

To investigate this, the experiments combined ERRANT-based grammatical evaluation with the Composite Over-Correction Index (OCI). This allowed the baseline, instruction-constrained, role-based, and few-shot prompts to be compared in terms of grammatical correction performance and utility-adjusted over-correction risk.

\section{Summary of Findings}

The results show that prompt wording had an effect on both correction accuracy and indicators of over-correction risk in this experimental setup.

The prompt refinement stage showed that even small changes in wording can affect model outputs. In the later comparison, the selected structured prompts performed better than the simple baseline. This suggests that the baseline prompt left the model with too much freedom to decide how much rewriting was appropriate.

The final strategy comparison showed clear differences between the prompt types. The baseline prompt produced the lowest $F_{0.5}$ score and the highest OCI, which suggests frequent but poorly targeted editing. The instruction-constrained prompt reduced OCI while still maintaining strong correction performance. The role-based prompt produced the best overall balance in this study, with high precision, strong $F_{0.5}$, low OCI, and the lowest average edit distance. The few-shot prompt achieved the highest $F_{0.5}$, but its higher OCI suggests that it encouraged a wider correction scope.

The sentence-length analysis showed that longer inputs were more challenging. Across the prompt conditions, long sentences produced lower $F_{0.5}$ scores and higher OCI values than short or medium sentences. This indicates that over-correction risk is influenced not only by prompt design, but also by the complexity of the input sentence.

The robustness analysis showed that the overall OCI ranking was not dependent on a single weighting scheme. The baseline remained the highest-risk condition, while the role-based and instruction-constrained prompts generally stayed among the lowest OCI conditions. This supports the use of OCI as a comparative indicator in this project.

The trade-off analysis showed that higher grammatical accuracy did not always require higher OCI. The role-based prompt achieved strong correction performance while keeping OCI low. By contrast, the baseline combined weak $F_{0.5}$ with the highest OCI, showing that making more changes does not necessarily lead to better correction.

The qualitative examples helped explain these numerical findings at sentence level. In some cases, the baseline paraphrased the learner's sentence, changed the style, or failed to follow the expected task format. These examples show why ERRANT alone was not sufficient for analysing correction behaviour in this project.

\section{Contributions}

This dissertation makes three main contributions.

First, it provides empirical evidence from the W\&I + LOCNESS subset used in this study that prompt wording can influence LLM-based GEC behaviour. Although the comparison is limited to one fixed OpenAI GPT-4o configuration, the differences between the baseline, instruction-constrained, role-based, and few-shot prompts were visible across ERRANT, OCI, and edit-distance results.

Second, the dissertation applies OCI as an indicator of over-correction risk. OCI combines source-output change with utility signals related to fluency improvement and meaning preservation. This made it possible to discuss rewriting behaviour alongside ERRANT accuracy, rather than treating $F_{0.5}$ as the only measure of success.

Third, the study offers a practical interpretation of the prompt strategies tested in the project. The baseline behaved as the least controlled condition. The instruction-constrained prompt was more conservative. The role-based prompt gave the strongest balance between accuracy and restraint. The few-shot prompt was the most accuracy-oriented, but it also appeared slightly more interventionist.

\section{Limitations}

This study has several limitations.

The first limitation is the use of a single dataset, W\&I + LOCNESS. Although this dataset is widely used in GEC research, the findings may not fully generalise to other learner corpora, writing levels, or domains.

The second limitation is that the experiments used one LLM configuration: OpenAI GPT-4o with fixed decoding settings. Other models may respond differently to the same prompts, particularly if they differ in instruction tuning or generation behaviour.

Another limitation is that OCI does not measure over-correction directly. It can indicate possible over-correction risk, but it cannot definitively judge whether each individual edit is necessary. Although the formula considers source-output change, fluency gain, and meaning preservation, it cannot always separate useful stylistic improvement from undesirable rewriting.

Finally, the study was carried out in a controlled evaluation setting. A real learner-facing GEC tool would involve additional factors, such as user interaction, how feedback is displayed, and more varied input quality.

\section{Future Work}

Several directions could extend this work.

One direction is to test the same prompt families across different models and datasets. This would help show whether the role-based and instruction-constrained prompts remain effective beyond the OpenAI GPT-4o configuration and W\&I + LOCNESS subset used in this dissertation.

Future work could also develop more detailed over-correction metrics. OCI is useful for comparing prompt conditions, but it cannot decide on its own whether a specific stylistic change is pedagogically helpful. Additional semantic or discourse-level features may help make this distinction clearer.

Another possible direction is to design adaptive prompts that respond to sentence length or complexity. The sentence-length analysis suggests that longer inputs may require stricter instructions about preserving the learner's wording.

User feedback would also make the evaluation more realistic. Learners or teachers could judge whether a correction is helpful, too minimal, or too interventionist. Automatic metrics are limited in how well they can represent this type of judgement.

\section{Final Remarks}

This dissertation shows that prompt engineering affects more than grammatical accuracy. In the experiments, different prompts also changed how much the model rewrote the learner's original sentence, and whether those changes were supported by fluency improvement and meaning preservation. Evaluating ERRANT performance together with OCI therefore gave a fuller view of correction behaviour than accuracy alone.

The main conclusion is that a useful learner-facing GEC prompt should not simply encourage the model to make more edits. It should guide the model to correct grammatical errors while preserving as much of the learner's own expression as possible.